\documentclass[10pt,journal,compsoc]{IEEEtran}

\usepackage[utf8]{inputenc}
\usepackage[T1]{fontenc}
\usepackage[spanish,es-tabla]{babel}
\usepackage{graphicx}
\usepackage{amsmath}
\usepackage{amssymb}
\usepackage{geometry}
\usepackage{hyperref}
\usepackage{booktabs}
\usepackage{color}

\geometry{margin=1in}

\begin{document}

\title{Manual de Operación de Usuario: Sistema conversacional y Dashboard de Detección de Fraude en WhatsApp}

\author{Esteban Mercado, Juan David Ocampo y Camilo}

\markboth{Manual de Usuario del Software,~Vol.~1, No.~1, Julio~2026}%
{Asistente de Detección de Fraude por WhatsApp}

\IEEEtitleabstractindextext{%
\begin{abstract}
Este manual describe el funcionamiento conversacional del bot de WhatsApp para la validación de identidad y detección de fraude en cédulas colombianas, así como la operación y visualización de indicadores en el panel de control web del proyecto. La guía detalla la interacción paso a paso del usuario final por chat y la explotación de métricas agregadas en tiempo real.
\end{abstract}
}

\maketitle

\IEEEdisplaynontitleabstractindextext
\IEEEpeerreviewmaketitle

\section{Descripción General}
\label{sec:overview}

El sistema **Detection Attacks** es un asistente automatizado accesible mediante la aplicación de mensajería **WhatsApp** cuyo propósito principal es guiar a los usuarios en el proceso de verificación y validación de su cédula colombiana para mitigar fraudes por suplantación de identidad.

\subsection{Características del Canal}
WhatsApp permite a los usuarios interactuar desde cualquier lugar del mundo con una experiencia móvil optimizada, sin fricciones de instalación de software propietario. El asistente virtual responde de forma interactiva y en tiempo real a las entradas de los usuarios, brindando respuestas amigables e intuitivas.

\subsection{Público Objetivo}
Este manual está dirigido a:
\begin{itemize}
    \item **Operadores del Sistema:** Personal encargado de monitorear el tráfico y revisar las métricas consolidadas de fraude.
    \item **Agentes de Soporte:** Personal que atiende solicitudes de usuarios finales cuyos documentos fueron rechazados por el bot.
\end{itemize}

\section{Guía de Interacción por WhatsApp}
\label{sec:interaction}

El bot opera como una conversación interactiva estructurada por estados. A continuación se detalla cómo operar el flujo conversacional.

\subsection{Paso 1: Saludo e Inicio del Proceso}
Para iniciar el flujo, el usuario puede enviar cualquier saludo común (como *"Hola"*, *"Buenas noches"*) o el comando específico **`reiniciar`**. El bot responderá saludándolo y solicitando la foto de la parte frontal de su cédula.

\subsection{Paso 2: Envío de la Cédula Frontal}
El usuario debe adjuntar y enviar la foto de la **cara frontal** de su cédula colombiana. 
\begin{itemize}
    \item **Si la foto es correcta:** El bot confirmará la recepción exitosa del documento frontal y solicitará que envíe la parte trasera.
    \item **Si hay problemas de calidad (score bajo) o errores en el análisis:** El bot le sugerirá amablemente tomar la foto nuevamente bajo mejores condiciones de iluminación y nitidez.
\end{itemize}

\subsection{Paso 3: Envío de la Cédula Trasera}
El usuario envía la foto de la **cara trasera** del documento. El sistema consolidará los análisis de ambas imágenes y emitirá un veredicto en tiempo real.

\subsection{Paso 4: Veredicto Final y Cierre}
- **Veredicto Exitoso (Auténtica):** El bot informará que el proceso se ha completado correctamente y se despedirá cordialmente.
- **Sospecha de Fraude o Rechazo:** Si alguna de las caras muestra indicadores de fraude (fotocopia a color, visualización en pantallas), el bot informará que la verificación automática no fue exitosa e instruirá al usuario sobre cómo contactar a soporte o reiniciar.

\subsection{Comando Especial: Reiniciar}
El usuario puede escribir la palabra **`reiniciar`** (no distingue mayúsculas) en cualquier etapa del proceso. Esto limpiará inmediatamente todos los datos y el progreso del usuario, restaurando su sesión al estado de saludo inicial.

\section{Guía del Dashboard de Monitoreo}
\label{sec:dashboard}

El Dashboard web de Detection Attacks proporciona visualizaciones en tiempo real del desempeño de las clasificaciones del bot. A continuación, se describen sus componentes principales.

\subsection{Indicadores Clave (KPIs)}
Ubicados en la parte superior del panel, muestran de forma resumida el estado del canal:
\begin{itemize}
    \item **Total Validations (Active Engine):** Conteo acumulado de transacciones recibidas.
    \item **Predicted Real / Authentic:** Cantidad de cédulas validadas como legítimas.
    \item **Predicted Fraud / Fake:** Cantidad de anomalías o ataques de suplantación detectados.
    \item **Percentage of Real (\%) / Percentage of Fraud (\%):** Proporción en porcentaje de cédulas aprobadas frente a sospechosas de fraude con precisión de 3 decimales.
\end{itemize}

\subsection{Histogramas de Distribución de Score}
Permiten analizar el comportamiento de las respuestas del modelo:
\begin{itemize}
    \item **Front/Reverse Prediction Score Distribution:** Dos gráficos de barra que exponen las predicciones en el rango continuo `[0.0, 1.0]`. 
    \item **Umbral $\tau = 0.5$:** Las barras a la izquierda del umbral representan documentos reales (color verde) y las que caen a la derecha indican intentos fraudulentos.
\end{itemize}

\subsection{Tendencia y Volumen Diario}
- **Daily Number of Fraud Predictions:** Gráfico de área apilado que muestra el volumen absoluto de validaciones aprobadas reales frente a fraudes detectados por día.
- **Daily Percentage of Fraud Detections:** Gráfico de área apilada al 100\% que facilita la visualización de la proporción diaria de fraudes y alertas detectadas sobre el volumen total.

\subsection{Monitoreo de Latencias y Tiempos de Respuesta}
El gráfico lineal de latencias grafica en milisegundos el tiempo consumido por los servicios:
\begin{itemize}
    \item **AWS Lambda (Rekognition):** Tiempo que toma procesar y clasificar la imagen.
    \item **Google Gemini:** Tiempo consumido para humanizar la respuesta final.
\end{itemize}

\subsection{Tabla de Historial (Summary of Document Validation)}
Muestra la bitácora consolidada de transacciones. Agrupa los eventos frontal y trasero de un usuario en un único registro para fácil auditoría. Se puede filtrar interactivamente por veredicto (*Todos*, *Real*, o *Fraude*).

\section{Solución de Problemas Comunes}
\label{sec:troubleshooting}

Esta sección resume las principales incidencias que pueden presentarse durante la operación del bot conversacional de WhatsApp y sus soluciones recomendadas.

\subsection{Problema 1: El bot no responde}
* **Causa:** El servidor de backend en FastAPI o la conexión de red está caída.
* **Solución:**
    1. Revisa el estado del servicio en Render (visita `https://detection-attacks-whatsapp.onrender.com/`).
    2. Verifica los logs de ejecución de FastAPI y asegúrate de que no haya un error de inicialización del pooler de la base de datos de Supabase.

\subsection{Problema 2: El bot rechaza fotos reales}
* **Causa:** Mala calidad en la captura de la cámara (borrosa, con reflejos de luz o inclinación pronunciada), lo que arroja un score erróneo por encima del umbral `0.5`.
* **Solución:**
    1. Aconseja al usuario tomar la foto con luz indirecta (evitando el flash directo sobre la cédula plástica).
    2. Mantener la cédula sobre un fondo oscuro y plano, alineada y enfocada en paralelo a la cámara.

\subsection{Problema 3: Error de Timeout de AWS Lambda (HTTP 500)}
* **Causa:** La Lambda `receiveCodedImage` en AWS se desborda en tiempo de ejecución (timeout inicial de 3 segundos superado).
* **Solución:**
    1. Asegúrate de configurar el timeout de la Lambda de AWS a 30 segundos y aumentar la memoria física a un mínimo de 512MB/1024MB en la consola de AWS. Esto evitará caídas por consumo de recursos al procesar imágenes de alta resolución.


\end{document}
